% Seminar 3 - Pagina 6. Fragment pentru
% % Seminar 3 - Pagina 6. Fragment pentru
% % Seminar 3 - Pagina 6. Fragment pentru
% % Seminar 3 - Pagina 6. Fragment pentru
% \input{seminars/pages/seminar3-pagina6-ro}; fără preambul sau pagină nouă.
% NOTE EDITORIALE:
% - Pagina continuă exemplele (i), (iii) și (iv) de pe pagina precedentă.
% - În prima linie, parantezele drepte și acoladele reprezintă partea
%   întreagă și partea fracționară; scrisul compact a fost interpretat
%   prin descompunerea x=[x]+{x}, y=[y]+{y}.
% - Manuscrisul verifică explicit reflexivitatea și simetria congruenței
%   modulo Z; tranzitivitatea nu este scrisă pe această pagină.
% - În margine se reamintește cazul (iv), x~y iff x-y in Q, și R/~ ~ R;
%   nu este repetat ca un exemplu numerotat separat.
% - Diagramele au fost redesenate pentru claritate; restricția lui p
%   la un sistem de reprezentanți este bijecția menționată.
% - Egalitatea argumentelor complexe cere o alegere consecventă a argumentului
%   (sau interpretarea modulo 2pi); notația manuscrisului nu precizează.
% - Eticheta aplicației către S' din ultima diagramă este palidă;
%   normalizarea z -> z/|z| corespunde construcției prezentate.

\subsection*{Seminar 3 - Pagina 6}
\begingroup
\setlength{\abovedisplayskip}{5pt}
\setlength{\belowdisplayskip}{5pt}
\setlength{\abovedisplayshortskip}{3pt}
\setlength{\belowdisplayshortskip}{3pt}

\noindent\textbf{Demonstrație (continuarea exemplului (i)).}
Scriem $x=\lfloor x\rfloor+\{x\}$ și $y=\lfloor y\rfloor+\{y\}$.
Dacă $x-y\in\mathbb Z$, atunci
\[
 \lfloor x\rfloor+\{x\}-\lfloor y\rfloor-\{y\}\in\mathbb Z,
 \qquad \{x\}-\{y\}\in\mathbb Z.
\]
Pentru relația $x\sim y\iff x-y\in\mathbb Z$:
\[
 \begin{aligned}
 \text{(R)}\quad &x\sim x &&(x-x=0\in\mathbb Z),\\
 \text{(S)}\quad &x\sim y\Longrightarrow y\sim x
     &&(x-y\in\mathbb Z\Longrightarrow y-x\in\mathbb Z).
 \end{aligned}
\]

\noindent\underline{Varianta 1.} Intervalul $[0,1)$ este sistem de
reprezentanți pentru $\sim$ dacă și numai dacă
\[
 p|_{[0,1)}:[0,1)\longrightarrow\mathbb R/{\sim}
 \quad\text{este bijecție}.
\]
\noindent\textbf{Pasul 1.} Pentru $x,y\in[0,1)$,
$x\ne y\Longrightarrow x\not\sim y$ (evident).

\noindent\textbf{Pasul 2.} Pentru orice $x\in\mathbb R$ există
$t\in[0,1)$ astfel încât $x-t\in\mathbb Z$: luăm $t=\{x\}$.

\medskip
\noindent\textbf{(iii) Numere complexe cu același modul.}
Pentru mulțimea factor $\mathbb C/{\sim}$, unde
$z_1\sim z_2\iff|z_1|=|z_2|$, considerăm
\[
 p|_{[0,+\infty)}:[0,+\infty)\longrightarrow\mathbb C/{\sim}.
\]
\noindent\textbf{(1)} Pentru $x,y\in[0,+\infty)$,
$x\ne y\Longrightarrow x\not\sim y$.

\noindent\textbf{(2)} Pentru orice $z\in\mathbb C$ există
$t\in[0,+\infty)$ astfel încât $z\sim t$: luăm $t=|z|$.
\[
\begin{tikzcd}[column sep=large,row sep=small]
 \mathbb C\arrow[r,"p"]\arrow[d,"{f(z)=|z|}"']
    & \mathbb C/{\sim} \\
 {[0,+\infty)}\arrow[ur,dashed,"{p|_{[0,+\infty)}}"'] &
\end{tikzcd}
\]
Dacă $|z_1|=|z_2|$, atunci $f(z_1)=f(z_2)$.

\medskip
\noindent\textbf{(iv) Numere complexe cu același argument.}
Notăm
\[
 S'=\{z\in\mathbb C\mid |z|=1\},\qquad
 z_1\sim z_2\iff\arg(z_1)=\arg(z_2)
 \quad(z_1,z_2\in\mathbb C^*).
\]
\noindent\textbf{(1)} Dacă $z_1,z_2\in S'$ și $z_1\ne z_2$,
atunci $z_1\not\sim z_2$ (argumentele sunt diferite).

\noindent\textbf{(2)} Pentru orice $z\in\mathbb C^*$ există
$z_1\in S'$ astfel încât $z\sim z_1$. Într-adevăr, scriem
\[
 z=r(\cos t+i\sin t),\qquad z_1=\cos t+i\sin t.
\]
Astfel, notițele identifică $\mathbb C^*/{\sim}$ cu $S'$ prin
reprezentanți de pe cercul unitate:
\[
\begin{tikzcd}[column sep=large,row sep=small]
 \mathbb C^*\arrow[r,"p"]\arrow[d,"{f(z)=z/|z|}"']
    & \mathbb C^*/{\sim} \\
 S'\arrow[ur,dashed,"{p|_{S'}}"'] &
\end{tikzcd}
\]

\endgroup
; fără preambul sau pagină nouă.
% NOTE EDITORIALE:
% - Pagina continuă exemplele (i), (iii) și (iv) de pe pagina precedentă.
% - În prima linie, parantezele drepte și acoladele reprezintă partea
%   întreagă și partea fracționară; scrisul compact a fost interpretat
%   prin descompunerea x=[x]+{x}, y=[y]+{y}.
% - Manuscrisul verifică explicit reflexivitatea și simetria congruenței
%   modulo Z; tranzitivitatea nu este scrisă pe această pagină.
% - În margine se reamintește cazul (iv), x~y iff x-y in Q, și R/~ ~ R;
%   nu este repetat ca un exemplu numerotat separat.
% - Diagramele au fost redesenate pentru claritate; restricția lui p
%   la un sistem de reprezentanți este bijecția menționată.
% - Egalitatea argumentelor complexe cere o alegere consecventă a argumentului
%   (sau interpretarea modulo 2pi); notația manuscrisului nu precizează.
% - Eticheta aplicației către S' din ultima diagramă este palidă;
%   normalizarea z -> z/|z| corespunde construcției prezentate.

\subsection*{Seminar 3 - Pagina 6}
\begingroup
\setlength{\abovedisplayskip}{5pt}
\setlength{\belowdisplayskip}{5pt}
\setlength{\abovedisplayshortskip}{3pt}
\setlength{\belowdisplayshortskip}{3pt}

\noindent\textbf{Demonstrație (continuarea exemplului (i)).}
Scriem $x=\lfloor x\rfloor+\{x\}$ și $y=\lfloor y\rfloor+\{y\}$.
Dacă $x-y\in\mathbb Z$, atunci
\[
 \lfloor x\rfloor+\{x\}-\lfloor y\rfloor-\{y\}\in\mathbb Z,
 \qquad \{x\}-\{y\}\in\mathbb Z.
\]
Pentru relația $x\sim y\iff x-y\in\mathbb Z$:
\[
 \begin{aligned}
 \text{(R)}\quad &x\sim x &&(x-x=0\in\mathbb Z),\\
 \text{(S)}\quad &x\sim y\Longrightarrow y\sim x
     &&(x-y\in\mathbb Z\Longrightarrow y-x\in\mathbb Z).
 \end{aligned}
\]

\noindent\underline{Varianta 1.} Intervalul $[0,1)$ este sistem de
reprezentanți pentru $\sim$ dacă și numai dacă
\[
 p|_{[0,1)}:[0,1)\longrightarrow\mathbb R/{\sim}
 \quad\text{este bijecție}.
\]
\noindent\textbf{Pasul 1.} Pentru $x,y\in[0,1)$,
$x\ne y\Longrightarrow x\not\sim y$ (evident).

\noindent\textbf{Pasul 2.} Pentru orice $x\in\mathbb R$ există
$t\in[0,1)$ astfel încât $x-t\in\mathbb Z$: luăm $t=\{x\}$.

\medskip
\noindent\textbf{(iii) Numere complexe cu același modul.}
Pentru mulțimea factor $\mathbb C/{\sim}$, unde
$z_1\sim z_2\iff|z_1|=|z_2|$, considerăm
\[
 p|_{[0,+\infty)}:[0,+\infty)\longrightarrow\mathbb C/{\sim}.
\]
\noindent\textbf{(1)} Pentru $x,y\in[0,+\infty)$,
$x\ne y\Longrightarrow x\not\sim y$.

\noindent\textbf{(2)} Pentru orice $z\in\mathbb C$ există
$t\in[0,+\infty)$ astfel încât $z\sim t$: luăm $t=|z|$.
\[
\begin{tikzcd}[column sep=large,row sep=small]
 \mathbb C\arrow[r,"p"]\arrow[d,"{f(z)=|z|}"']
    & \mathbb C/{\sim} \\
 {[0,+\infty)}\arrow[ur,dashed,"{p|_{[0,+\infty)}}"'] &
\end{tikzcd}
\]
Dacă $|z_1|=|z_2|$, atunci $f(z_1)=f(z_2)$.

\medskip
\noindent\textbf{(iv) Numere complexe cu același argument.}
Notăm
\[
 S'=\{z\in\mathbb C\mid |z|=1\},\qquad
 z_1\sim z_2\iff\arg(z_1)=\arg(z_2)
 \quad(z_1,z_2\in\mathbb C^*).
\]
\noindent\textbf{(1)} Dacă $z_1,z_2\in S'$ și $z_1\ne z_2$,
atunci $z_1\not\sim z_2$ (argumentele sunt diferite).

\noindent\textbf{(2)} Pentru orice $z\in\mathbb C^*$ există
$z_1\in S'$ astfel încât $z\sim z_1$. Într-adevăr, scriem
\[
 z=r(\cos t+i\sin t),\qquad z_1=\cos t+i\sin t.
\]
Astfel, notițele identifică $\mathbb C^*/{\sim}$ cu $S'$ prin
reprezentanți de pe cercul unitate:
\[
\begin{tikzcd}[column sep=large,row sep=small]
 \mathbb C^*\arrow[r,"p"]\arrow[d,"{f(z)=z/|z|}"']
    & \mathbb C^*/{\sim} \\
 S'\arrow[ur,dashed,"{p|_{S'}}"'] &
\end{tikzcd}
\]

\endgroup
; fără preambul sau pagină nouă.
% NOTE EDITORIALE:
% - Pagina continuă exemplele (i), (iii) și (iv) de pe pagina precedentă.
% - În prima linie, parantezele drepte și acoladele reprezintă partea
%   întreagă și partea fracționară; scrisul compact a fost interpretat
%   prin descompunerea x=[x]+{x}, y=[y]+{y}.
% - Manuscrisul verifică explicit reflexivitatea și simetria congruenței
%   modulo Z; tranzitivitatea nu este scrisă pe această pagină.
% - În margine se reamintește cazul (iv), x~y iff x-y in Q, și R/~ ~ R;
%   nu este repetat ca un exemplu numerotat separat.
% - Diagramele au fost redesenate pentru claritate; restricția lui p
%   la un sistem de reprezentanți este bijecția menționată.
% - Egalitatea argumentelor complexe cere o alegere consecventă a argumentului
%   (sau interpretarea modulo 2pi); notația manuscrisului nu precizează.
% - Eticheta aplicației către S' din ultima diagramă este palidă;
%   normalizarea z -> z/|z| corespunde construcției prezentate.

\subsection*{Seminar 3 - Pagina 6}
\begingroup
\setlength{\abovedisplayskip}{5pt}
\setlength{\belowdisplayskip}{5pt}
\setlength{\abovedisplayshortskip}{3pt}
\setlength{\belowdisplayshortskip}{3pt}

\noindent\textbf{Demonstrație (continuarea exemplului (i)).}
Scriem $x=\lfloor x\rfloor+\{x\}$ și $y=\lfloor y\rfloor+\{y\}$.
Dacă $x-y\in\mathbb Z$, atunci
\[
 \lfloor x\rfloor+\{x\}-\lfloor y\rfloor-\{y\}\in\mathbb Z,
 \qquad \{x\}-\{y\}\in\mathbb Z.
\]
Pentru relația $x\sim y\iff x-y\in\mathbb Z$:
\[
 \begin{aligned}
 \text{(R)}\quad &x\sim x &&(x-x=0\in\mathbb Z),\\
 \text{(S)}\quad &x\sim y\Longrightarrow y\sim x
     &&(x-y\in\mathbb Z\Longrightarrow y-x\in\mathbb Z).
 \end{aligned}
\]

\noindent\underline{Varianta 1.} Intervalul $[0,1)$ este sistem de
reprezentanți pentru $\sim$ dacă și numai dacă
\[
 p|_{[0,1)}:[0,1)\longrightarrow\mathbb R/{\sim}
 \quad\text{este bijecție}.
\]
\noindent\textbf{Pasul 1.} Pentru $x,y\in[0,1)$,
$x\ne y\Longrightarrow x\not\sim y$ (evident).

\noindent\textbf{Pasul 2.} Pentru orice $x\in\mathbb R$ există
$t\in[0,1)$ astfel încât $x-t\in\mathbb Z$: luăm $t=\{x\}$.

\medskip
\noindent\textbf{(iii) Numere complexe cu același modul.}
Pentru mulțimea factor $\mathbb C/{\sim}$, unde
$z_1\sim z_2\iff|z_1|=|z_2|$, considerăm
\[
 p|_{[0,+\infty)}:[0,+\infty)\longrightarrow\mathbb C/{\sim}.
\]
\noindent\textbf{(1)} Pentru $x,y\in[0,+\infty)$,
$x\ne y\Longrightarrow x\not\sim y$.

\noindent\textbf{(2)} Pentru orice $z\in\mathbb C$ există
$t\in[0,+\infty)$ astfel încât $z\sim t$: luăm $t=|z|$.
\[
\begin{tikzcd}[column sep=large,row sep=small]
 \mathbb C\arrow[r,"p"]\arrow[d,"{f(z)=|z|}"']
    & \mathbb C/{\sim} \\
 {[0,+\infty)}\arrow[ur,dashed,"{p|_{[0,+\infty)}}"'] &
\end{tikzcd}
\]
Dacă $|z_1|=|z_2|$, atunci $f(z_1)=f(z_2)$.

\medskip
\noindent\textbf{(iv) Numere complexe cu același argument.}
Notăm
\[
 S'=\{z\in\mathbb C\mid |z|=1\},\qquad
 z_1\sim z_2\iff\arg(z_1)=\arg(z_2)
 \quad(z_1,z_2\in\mathbb C^*).
\]
\noindent\textbf{(1)} Dacă $z_1,z_2\in S'$ și $z_1\ne z_2$,
atunci $z_1\not\sim z_2$ (argumentele sunt diferite).

\noindent\textbf{(2)} Pentru orice $z\in\mathbb C^*$ există
$z_1\in S'$ astfel încât $z\sim z_1$. Într-adevăr, scriem
\[
 z=r(\cos t+i\sin t),\qquad z_1=\cos t+i\sin t.
\]
Astfel, notițele identifică $\mathbb C^*/{\sim}$ cu $S'$ prin
reprezentanți de pe cercul unitate:
\[
\begin{tikzcd}[column sep=large,row sep=small]
 \mathbb C^*\arrow[r,"p"]\arrow[d,"{f(z)=z/|z|}"']
    & \mathbb C^*/{\sim} \\
 S'\arrow[ur,dashed,"{p|_{S'}}"'] &
\end{tikzcd}
\]

\endgroup
; fără preambul sau pagină nouă.
% NOTE EDITORIALE:
% - Pagina continuă exemplele (i), (iii) și (iv) de pe pagina precedentă.
% - În prima linie, parantezele drepte și acoladele reprezintă partea
%   întreagă și partea fracționară; scrisul compact a fost interpretat
%   prin descompunerea x=[x]+{x}, y=[y]+{y}.
% - Manuscrisul verifică explicit reflexivitatea și simetria congruenței
%   modulo Z; tranzitivitatea nu este scrisă pe această pagină.
% - În margine se reamintește cazul (iv), x~y iff x-y in Q, și R/~ ~ R;
%   nu este repetat ca un exemplu numerotat separat.
% - Diagramele au fost redesenate pentru claritate; restricția lui p
%   la un sistem de reprezentanți este bijecția menționată.
% - Egalitatea argumentelor complexe cere o alegere consecventă a argumentului
%   (sau interpretarea modulo 2pi); notația manuscrisului nu precizează.
% - Eticheta aplicației către S' din ultima diagramă este palidă;
%   normalizarea z -> z/|z| corespunde construcției prezentate.

\subsection*{Seminar 3 - Pagina 6}
\begingroup
\setlength{\abovedisplayskip}{5pt}
\setlength{\belowdisplayskip}{5pt}
\setlength{\abovedisplayshortskip}{3pt}
\setlength{\belowdisplayshortskip}{3pt}

\noindent\textbf{Demonstrație (continuarea exemplului (i)).}
Scriem $x=\lfloor x\rfloor+\{x\}$ și $y=\lfloor y\rfloor+\{y\}$.
Dacă $x-y\in\mathbb Z$, atunci
\[
 \lfloor x\rfloor+\{x\}-\lfloor y\rfloor-\{y\}\in\mathbb Z,
 \qquad \{x\}-\{y\}\in\mathbb Z.
\]
Pentru relația $x\sim y\iff x-y\in\mathbb Z$:
\[
 \begin{aligned}
 \text{(R)}\quad &x\sim x &&(x-x=0\in\mathbb Z),\\
 \text{(S)}\quad &x\sim y\Longrightarrow y\sim x
     &&(x-y\in\mathbb Z\Longrightarrow y-x\in\mathbb Z).
 \end{aligned}
\]

\noindent\underline{Varianta 1.} Intervalul $[0,1)$ este sistem de
reprezentanți pentru $\sim$ dacă și numai dacă
\[
 p|_{[0,1)}:[0,1)\longrightarrow\mathbb R/{\sim}
 \quad\text{este bijecție}.
\]
\noindent\textbf{Pasul 1.} Pentru $x,y\in[0,1)$,
$x\ne y\Longrightarrow x\not\sim y$ (evident).

\noindent\textbf{Pasul 2.} Pentru orice $x\in\mathbb R$ există
$t\in[0,1)$ astfel încât $x-t\in\mathbb Z$: luăm $t=\{x\}$.

\medskip
\noindent\textbf{(iii) Numere complexe cu același modul.}
Pentru mulțimea factor $\mathbb C/{\sim}$, unde
$z_1\sim z_2\iff|z_1|=|z_2|$, considerăm
\[
 p|_{[0,+\infty)}:[0,+\infty)\longrightarrow\mathbb C/{\sim}.
\]
\noindent\textbf{(1)} Pentru $x,y\in[0,+\infty)$,
$x\ne y\Longrightarrow x\not\sim y$.

\noindent\textbf{(2)} Pentru orice $z\in\mathbb C$ există
$t\in[0,+\infty)$ astfel încât $z\sim t$: luăm $t=|z|$.
\[
\begin{tikzcd}[column sep=large,row sep=small]
 \mathbb C\arrow[r,"p"]\arrow[d,"{f(z)=|z|}"']
    & \mathbb C/{\sim} \\
 {[0,+\infty)}\arrow[ur,dashed,"{p|_{[0,+\infty)}}"'] &
\end{tikzcd}
\]
Dacă $|z_1|=|z_2|$, atunci $f(z_1)=f(z_2)$.

\medskip
\noindent\textbf{(iv) Numere complexe cu același argument.}
Notăm
\[
 S'=\{z\in\mathbb C\mid |z|=1\},\qquad
 z_1\sim z_2\iff\arg(z_1)=\arg(z_2)
 \quad(z_1,z_2\in\mathbb C^*).
\]
\noindent\textbf{(1)} Dacă $z_1,z_2\in S'$ și $z_1\ne z_2$,
atunci $z_1\not\sim z_2$ (argumentele sunt diferite).

\noindent\textbf{(2)} Pentru orice $z\in\mathbb C^*$ există
$z_1\in S'$ astfel încât $z\sim z_1$. Într-adevăr, scriem
\[
 z=r(\cos t+i\sin t),\qquad z_1=\cos t+i\sin t.
\]
Astfel, notițele identifică $\mathbb C^*/{\sim}$ cu $S'$ prin
reprezentanți de pe cercul unitate:
\[
\begin{tikzcd}[column sep=large,row sep=small]
 \mathbb C^*\arrow[r,"p"]\arrow[d,"{f(z)=z/|z|}"']
    & \mathbb C^*/{\sim} \\
 S'\arrow[ur,dashed,"{p|_{S'}}"'] &
\end{tikzcd}
\]

\endgroup
